\begin{explanationbox}
	\textbf{\textcolor{CExplanation}{一句话直觉}}
	通径是圆锥曲线中结构最简单的焦点弦，其长度可直接由曲线基本参数（$a,b$ 或 $p$）算出，无需额外联立消元。

	\medskip
	\textbf{\textcolor{CExplanation}{核心拆解}}
	\begin{description}[style=nextline, leftmargin=2.8em, labelsep=0.5em, itemsep=0.45em]
		\item[\textbf{要点一（过焦点条件）：}]
			通径必须经过焦点，即弦所在直线通过曲线的某一焦点。
		\item[\textbf{要点二（垂直条件）：}]
			通径与焦点所在坐标轴垂直，因此横坐标（或纵坐标）固定为焦点坐标。
	\end{description}

	\medskip
	\textbf{\textcolor{CExplanation}{几何本质（垂直简化）}}
	当弦垂直于坐标轴时，弦两端点的横坐标相同，代入曲线方程后可直接求解纵坐标，大大简化了计算。

	\medskip
	\textbf{\textcolor{CExplanation}{代数意义（代入消元）}}
	将焦点横坐标（如 $x=c$ 或 $x=p/2$）代入标准方程，转化为关于纵坐标的一元二次方程，直接解出 $y$，再由 $|y_1-y_2|$ 得弦长。

	\medskip
	\textbf{\textcolor{CExplanation}{考点价值（快速计算）}}
	\begin{itemize}[leftmargin=*, itemsep=0.5em, topsep=0.4em]
		\item \textbf{考法一（直接求通径）：} 给出标准方程，直接套用 $L=2b^2/a$ 或 $L=2p$ 计算通径长。
		\item \textbf{考法二（方程求解参数）：} 利用通径长列方程，求解曲线中的未知参数（如 $a,b,p$）。
	\end{itemize}

	\medskip
	\textbf{\textcolor{CExplanation}{顿悟点}}
	通径长公式中不显含焦距 $c$，椭圆和双曲线的通径长只由 $a,b$ 决定，抛物线的通径长仅由 $p$ 决定。这意味着通径长度只与曲线的“开口形状”有关，而与焦点具体位置无关（对于同一曲线）。

	\medskip
	\textbf{\textcolor{CExplanation}{使用场景}}
	\begin{itemize}[leftmargin=*, itemsep=0.5em, topsep=0.4em]
		\item \textbf{场景一（标准曲线套用）：} 遇到标准方程且已知弦过焦点并垂直坐标轴时，直接应用通径公式。
		\item \textbf{场景二（焦点弦定位）：} 当题目描述“过焦点且垂直于对称轴的弦”时，优先使用通径长度结论。
	\end{itemize}

\end{explanationbox}
